%
\newacronym{mat}{Matlab}{Matrix Laboratory}
\newacronym{ide}{IDE}{Integrated Development Environment}

% === Custom Chapter / Module Command ===
\newcounter{chapt}
\renewcommand{\thechapt}{\Roman{chapt}} % Usa numeri romani (I, II, III). Cambia in \arabic per (1, 2, 3)

% Usiamo \NewDocumentCommand da xparse: 'o' = optional, 'm' = mandatory
\NewDocumentCommand{\chapt}{ o m }{%
    \clearpage % Forza l'inizio su una nuova pagina
    \thispagestyle{empty} % Rimuove header e footer da questa pagina di "copertina"
    \refstepcounter{chapt}

    \vspace*{0.15\textheight} % Spazio verticale per centrare otticamente il titolo

    \begin{flushleft}
        {\xcharter\huge\bfseries \textcolor{black!60}{Modulo \thechapt}} \par\vspace{0.5cm}
        {\sffamily\Huge\bfseries #2} \par

        % \IfValueT controlla se il parametro opzionale 'o' è stato inserito
        \IfValueT{#1}{%
            \vspace{0.3cm}
            {\xcharter\Large\itshape Prof. #1} % Usa \sffamily per chiamare il Source Sans Pro
        }
    \end{flushleft}

    \vspace{2cm}

    % Aggiunge la voce all'indice (TOC).
    \addcontentsline{toc}{part}{Modulo \thechapt: #2}

    % --- Gestione Contatori ---
    \setcounter{lecture}{0}    % Fa ripartire le lezioni da 1 per ogni nuovo professore
    \setcounter{assignment}{0} % Fa ripartire gli assignment da 1
    \setcounter{section}{0}    % Fa ripartire le sezioni da 1

    \vfill
    \clearpage % Forza l'inizio del testo vero e proprio sulla pagina successiva
}